\documentclass[12pt]{extarticle}

% ===== Encodage & langue =====
\usepackage{fontspec}
\usepackage{polyglossia}
\setdefaultlanguage{french}

% ===== Police Verdana (compiler avec XeLaTeX ou LuaLaTeX) =====
\setsansfont{Verdana}
\renewcommand{\familydefault}{\sfdefault}

% ===== Interligne =====
\usepackage{setspace}
\setstretch{1.1}

% ===== Marges réduites =====
\usepackage{geometry}
\geometry{a4paper, top=1.8cm, bottom=1.8cm, left=2cm, right=2cm}

% ===== Mathématiques =====
\usepackage{amsmath}

% ===== Espacement réduit autour des équations =====
\setlength{\abovedisplayskip}{4pt}
\setlength{\belowdisplayskip}{4pt}
\setlength{\abovedisplayshortskip}{2pt}
\setlength{\belowdisplayshortskip}{2pt}

% ===== Espacement réduit des sections =====
\usepackage{titlesec}
\titlespacing*{\section}{0pt}{1.2ex}{0.6ex}

% ===== Titre du document =====
\title{\textbf{Correction du contrôle séquence 1 version 7}}
\author{}
\date{}

\begin{document}
\maketitle

% --------------------------------------------------
\section*{Exercice 1}
\begin{align*}
13 - 6 - 3 + 28 - 4 &= 7 - 3 + 28 - 4 \\
                     &= 4 + 28 - 4 \\
                     &= 32 - 4 \\
                     &= 28
\end{align*}

% --------------------------------------------------
\section*{Exercice 2}

\begin{align*}
\textbf{a)}
5 + 6 \times 4 &= 5 + 24 \\
               &= 29
\end{align*}

\textbf{b)}
\begin{align*}
43 - 8 \times 5 &= 43 - 40 \\
                &= 3
\end{align*}

\textbf{c)}
\begin{align*}
45 + 24 \div 8 &= 45 + 3 \\
               &= 48
\end{align*}

\textbf{d)}
\begin{align*}
5 + 7 \times 6 + 8 - 6 \div 2 &= 5 + 42 + 8 - 3 \\
                              &= 52
\end{align*}

% --------------------------------------------------
\section*{Exercice 3}
\begin{align*}
4 \times 7 + 3 \times 8 &= 28 + 24 \\
                        &= 52
\end{align*}

La dépense est de 52~€.

% --------------------------------------------------
\section*{Exercice 4}

\textbf{a)} $33 - (6 + 8) \div 2$
\begin{align*}
&= 33 - 14 \div 2 \\
&= 33 - 7 \\
&= 26
\end{align*}

\textbf{b)} $(4 + 9) \times 5 - 6 + 8$
\begin{align*}
&= 13 \times 5 - 6 + 8 \\
&= 65 - 6 + 8 \\
&= 67
\end{align*}

\textbf{c)} $(26 - (7 + 8)) \times 5$
\begin{align*}
&= (26 - 15) \times 5 \\
&= 11 \times 5 \\
&= 55
\end{align*}

\textbf{d)} $5 \times (23 - (6 + 11)) - 8$
\begin{align*}
&= 5 \times (23 - 17) - 8 \\
&= 5 \times 6 - 8 \\
&= 30 - 8 \\
&= 22
\end{align*}

% --------------------------------------------------
\section*{Exercice 5}
\begin{align*}
14 - 4 \times 3 &= 14 - 12 \\
                &= 2
\end{align*}

Une gomme coûte 2~€.

% --------------------------------------------------
\section*{Exercice bonus 1}

\textbf{a)} $9 \times 4 + 16$

\textbf{b)} $8 \times (7 + 5)$

% --------------------------------------------------
\section*{Exercice bonus 2}

\textbf{a)} La somme du produit de 9 par 5 et de 7.

\textbf{b)} Le produit de la somme de 6 et de 9 par 8.

\end{document}